\documentclass[tikz,border=8pt]{standalone}
\usepackage{tikz}
\usepackage{amsmath}
\usepackage{amssymb}
\usepackage{ctex}
\usetikzlibrary{calc, positioning, arrows.meta, decorations.pathreplacing, fit, backgrounds}

\begin{document}
\begin{tikzpicture}[>=Stealth]

%% ===== 上横轴：R^n（n=5）=====
% 总长度 5 单位，行空间 r=3，零空间 n-r=2

\def\totaln{5}
\def\rankr{3}
\def\bary{2.5}    % 横轴 y 坐标
\def\barh{0.7}    % 条形高度
\def\xscale{1.4}  % 每单位的像素宽度（cm）

% 行空间段（0 到 r=3），蓝色
\fill[blue!30] (0, \bary) rectangle ({\rankr*\xscale}, {\bary+\barh});
\draw[blue!60, thick] (0, \bary) rectangle ({\rankr*\xscale}, {\bary+\barh});

% 零空间段（r=3 到 n=5），橙色
\fill[orange!30] ({\rankr*\xscale}, \bary) rectangle ({\totaln*\xscale}, {\bary+\barh});
\draw[orange!60, thick] ({\rankr*\xscale}, \bary) rectangle ({\totaln*\xscale}, {\bary+\barh});

% 标注：行空间（主标签在条形内，维数标注在上方）
\node[font=\small\bfseries, blue!80!black] at ({(\rankr*\xscale)/2}, {\bary+\barh*0.65})
  {行空间 $C(A^\top)$};
\node[font=\scriptsize, blue!70!black] at ({(\rankr*\xscale)/2}, {\bary+\barh*0.2})
  {$\mathrm{rank}\;r = 3$};

% 标注：零空间（主标签在条形内，维数标注在条形内）
\node[font=\small\bfseries, orange!80!black]
  at ({(\rankr+(\totaln-\rankr)/2)*\xscale}, {\bary+\barh*0.65})
  {零空间 $N(A)$};
\node[font=\scriptsize, orange!70!black]
  at ({(\rankr+(\totaln-\rankr)/2)*\xscale}, {\bary+\barh*0.2})
  {$n-r = 2$};

% 总宽度标注（大括号）
\draw[decorate, decoration={brace, amplitude=6pt}, gray!60]
  (0, {\bary+\barh+0.1}) -- ({\totaln*\xscale}, {\bary+\barh+0.1})
  node[midway, above=6pt, font=\small, gray!80] {$\mathbb{R}^n$（$n=5$）};

% 分隔线标注
\draw[gray!50, dashed, thin] ({\rankr*\xscale}, {\bary-0.5}) -- ({\rankr*\xscale}, {\bary+\barh+0.7});
\node[font=\scriptsize, gray!70] at ({\rankr*\xscale}, {\bary+\barh+0.85}) {$r=3$};

% 轴线
\draw[->, gray!60, thin] (-0.3, \bary) -- ({\totaln*\xscale+0.5}, \bary)
  node[right, font=\scriptsize, gray] {维度轴};
\foreach \i in {0,1,2,3,4,5}{
  \draw[gray!50, thin] ({\i*\xscale}, {\bary-0.08}) -- ({\i*\xscale}, {\bary+0.08});
  \node[font=\tiny, gray!60] at ({\i*\xscale}, {\bary-0.2}) {\i};
}

%% ===== 下横轴：R^m（m=4）=====
\def\totalm{4}
\def\nullitym{1}   % m-r = 4-3 = 1
\def\baryb{-0.5}  % 下轴 y 坐标

% 列空间段（0 到 r=3），绿色
\fill[teal!30] (0, \baryb) rectangle ({\rankr*\xscale}, {\baryb+\barh});
\draw[teal!60, thick] (0, \baryb) rectangle ({\rankr*\xscale}, {\baryb+\barh});

% 左零空间段（r=3 到 m=4），红色
\fill[red!20] ({\rankr*\xscale}, \baryb) rectangle ({\totalm*\xscale}, {\baryb+\barh});
\draw[red!50, thick] ({\rankr*\xscale}, \baryb) rectangle ({\totalm*\xscale}, {\baryb+\barh});

% 标注：列空间（主标签 + 维数都放条形内）
\node[font=\small\bfseries, teal!80!black] at ({(\rankr*\xscale)/2}, {\baryb+\barh*0.65})
  {列空间 $C(A)$};
\node[font=\scriptsize, teal!70!black] at ({(\rankr*\xscale)/2}, {\baryb+\barh*0.2})
  {$\mathrm{rank}\;r = 3$};

% 标注：左零空间
\node[font=\small\bfseries, red!70!black]
  at ({(\rankr+0.5)*\xscale}, {\baryb+\barh*0.65})
  {左零空间};
\node[font=\scriptsize, red!60]
  at ({(\rankr+0.5)*\xscale}, {\baryb+\barh*0.2})
  {$m-r = 1$};

% 总宽度标注（大括号）
\draw[decorate, decoration={brace, amplitude=6pt, mirror}, gray!60]
  (0, {\baryb-0.1}) -- ({\totalm*\xscale}, {\baryb-0.1})
  node[midway, below=6pt, font=\small, gray!80] {$\mathbb{R}^m$（$m=4$）};

% 分隔线标注
\draw[gray!50, dashed, thin] ({\rankr*\xscale}, {\baryb-0.7}) -- ({\rankr*\xscale}, {\baryb+\barh+0.5});

% 轴线
\draw[->, gray!60, thin] (-0.3, {\baryb+\barh}) -- ({\totalm*\xscale+0.5}, {\baryb+\barh})
  node[right, font=\scriptsize, gray] {维度轴};
\foreach \i in {0,1,2,3,4}{
  \draw[gray!50, thin] ({\i*\xscale}, {\baryb+\barh-0.08}) -- ({\i*\xscale}, {\baryb+\barh+0.08});
  \node[font=\tiny, gray!60] at ({\i*\xscale}, {\baryb+\barh+0.18}) {\i};
}

%% ===== 连线：同维度对应关系 =====

% 行空间 ↔ 列空间（维数同为 r=3）
\draw[<->, blue!60, thick, dashed]
  ({(\rankr*\xscale)/2}, \bary-0.05) -- ({(\rankr*\xscale)/2}, {\baryb+\barh+0.05})
  node[midway, right=4pt, font=\scriptsize, blue!70!black]
  {$A$ 的满射：行空间$\xrightarrow{\sim}$列空间};

% 零空间 ↔ 左零空间（维数差异：2 vs 1，仅标注对应关系）
\draw[<->, orange!60, thick, dashed]
  ({(\rankr+(\totaln-\rankr)/2)*\xscale}, \bary-0.05)
  to[bend right=20]
  node[midway, right=4pt, font=\scriptsize, orange!80!black]
  {补空间（均映到 $\mathbf{0}$）}
  ({(\rankr+0.5)*\xscale}, {\baryb+\barh+0.05});

%% ===== 定理框 =====
\node[font=\small, fill=white, draw=gray!50, thin, rounded corners=3pt, inner sep=8pt,
      align=center] at ({2.5*\xscale}, -2.5) {
  \textbf{秩-零化度定理（Rank-Nullity Theorem）}\\[3pt]
  $\dim C(A^\top) + \dim N(A) = n$
  \quad $\Longrightarrow$ \quad $3 + 2 = 5$ \checkmark\\[2pt]
  $\dim C(A) + \dim N(A^\top) = m$
  \quad $\Longrightarrow$ \quad $3 + 1 = 4$ \checkmark
};

%% 整体标题
\node[font=\large\bfseries] at ({2.5*\xscale}, 4.5)
  {秩-零化度定理的维度分解（$m=4,\,n=5,\,r=3$）};

\end{tikzpicture}
\end{document}
